% architecture.tex — class overview and interaction diagrams for ComplexGitSync

\chapter{Architecture Overview}
\label{chap:architecture}

This chapter provides a logical scheme of the main classes in \cgspkg{} and
explains how they interact.  Each domain class lives in its own
\texttt{.py} module, following the
\emph{one class per source file} convention of \texttt{DevSpecs.md}.

% -----------------------------------------------------------------------
\section{Module Layout}

The source tree under \texttt{src/ComplexGitSync/} maps directly to the class
hierarchy: one \texttt{.py} file per domain class, with thin re-export shims
(\texttt{documents.py}, \texttt{registry.py}) retained for backward
compatibility.

\begin{lstlisting}[language=bash,caption={Source tree — one class per module}]
src/ComplexGitSync/
  # Entry points
  __init__.py              # public API surface
  cli.py                   # cgitsync CLI (build_parser / main)
  client.py                # ComplexGitSyncClient

  # Core model — centred on GitTree
  git_tree.py              # GitTree
  git_repo.py              # GitRepo
  orchestre.py             # Orchestre
  git_runner.py            # GitRunner
  repo_address.py          # RepoAddress

  # Enumeration classes
  access_protocol.py       # AccessProtocol
  git_provider.py          # GitProvider
  tree_lifecycle_state.py  # TreeLifecycleState
  repo_lifecycle_state.py  # RepoLifecycleState
  discovery_state.py       # DiscoveryState
  sync_state.py            # SyncState
  ref_kind.py              # RefKind
  node_type.py             # NodeType

  # Registry model
  repo_node.py             # RepoNode
  repo_registry_entry.py   # RepoRegistryEntry
  dependency_tree_registry.py  # DependencyTreeRegistry
  project_tree_state.py    # ProjectTreeState
  registry.py              # builder functions + re-export shim

  # Document classes
  config_document.py       # ConfigDocument (base)
  cgs_document.py          # CgsDocument  (.cgs authoring spec)
  gts_document.py          # GtsDocument  (.gts snapshot)
  goc_document.py          # GocDocument  (.goc command script)
  documents.py             # re-export shim

  # Infrastructure
  run_logger.py            # CommandRunLogger + create_run_logger
  state_store.py           # RuntimeStateStore
  discovery.py             # discover_nested_configs()
  render.py                # format_project_tree / format_registry_json
  errors.py                # exception hierarchy
\end{lstlisting}

% -----------------------------------------------------------------------
\section{Core Object Model}

Everything in \cgspkg{} is centred on \textbf{GitTree}, which represents the
dependency graph of all repositories in the project.
Figure~\ref{fig:core-model} shows the full ownership chain from the top-level
client down to individual repository objects.

\begin{itemize}
  \item \texttt{ComplexGitSyncClient} is the public facade.  All CLI commands
    delegate to its methods.
  \item \texttt{Orchestre} is the coordination layer.  It owns exactly one
    \texttt{GitTree} and exposes repository registration to the client.
  \item \texttt{GitTree} holds a dictionary of \texttt{GitRepo} objects, keyed
    by \texttt{project\_name}.  It owns the lifecycle of in-memory repo
    identity and supports correction operations (\texttt{force\_repo\_sha},
    \texttt{force\_repo\_keys}).
  \item \texttt{GitRepo} is a dataclass capturing the static identity of a
    single repository: provider, namespace, project name, access protocol, and
    commit SHA.
  \item \texttt{RepoAddress} derives the remote URL (SSH or HTTPS) from a
    \texttt{GitRepo} without mutating it.
\end{itemize}

% class_diagram.tex  — TikZ source for the ComplexGitSync class overview figure.
%
% This file is \input{} from docs/architecture.tex.  It draws three related
% diagrams on separate pages so each fits comfortably within an A4 margin:
%
%   Figure 1 — Core object model centred on GitTree
%   Figure 2 — Document hierarchy (ConfigDocument and subclasses)
%   Figure 3 — Registry model (WorkingGitTree and supporting types)

% ============================================================
%  Shared TikZ style definitions
% ============================================================

\tikzset{
  %— UML-style class box with three compartments
  classbox/.style={
    rectangle split,
    rectangle split parts=3,
    rectangle split part fill={blue!10,white,white},
    draw=blue!40,
    thick,
    font=\ttfamily\footnotesize,
    text width=4.8cm,
    align=left,
    inner sep=4pt,
  },
  enumbox/.style={
    rectangle split,
    rectangle split parts=2,
    rectangle split part fill={orange!15,white},
    draw=orange!50,
    thick,
    font=\ttfamily\footnotesize,
    text width=3.8cm,
    align=left,
    inner sep=4pt,
  },
  %— Arrows
  hasarrow/.style={-{Stealth[length=6pt]}, thick},
  inharrow/.style={-{Triangle[open,length=7pt]}, thick},
  comporr/.style={-{Diamond[length=7pt]}, thick},
}

% ============================================================
%  Figure 1 — Core object model centred on GitTree
% ============================================================

\begin{figure}[ht]
\centering
\caption{Core object model: \texttt{ComplexGitSyncClient} orchestrates through
  \texttt{Orchestre}, which owns the central \texttt{GitTree} composed of
  \texttt{GitRepo} instances.  \texttt{RepoAddress} derives the remote URL from
  a \texttt{GitRepo}.}
\label{fig:core-model}

\begin{tikzpicture}[node distance=1.2cm and 2.2cm]

% --- ComplexGitSyncClient ---
\node[classbox] (client) {
  \textbf{ComplexGitSyncClient}
  \nodepart{two}
  orchestre: Orchestre\\
  git\_runner: GitRunner\\
  state\_store: RuntimeStateStore\\
  registry: WorkingGitTree\\
  run\_logger: CommandRunLogger
  \nodepart{three}
  read(cgs) / load\_gts(gts)\\
  initialise(cgs|gts) / load(cgs|gts)\\
  pull(cgs|gts) / clone(cgs)\\
  checkout(ref) / add()\\
  git(tree, cmd, args)\\
  freeze(name)
};

% --- Orchestre ---
\node[classbox, below=1.8cm of client] (orchestre) {
  \textbf{Orchestre}
  \nodepart{two}
  git\_tree: GitTree
  \nodepart{three}
  register\_repo(repo)
};

% --- GitTree ---
\node[classbox, below=1.8cm of orchestre] (gittree) {
  \textbf{GitTree}
  \nodepart{two}
  repos: dict[str, GitRepo]
  \nodepart{three}
  add\_repo(repo)\\
  force\_repo\_sha(name, sha)\\
  force\_repo\_keys(name, **)
};

% --- GitRepo ---
\node[classbox, right=2.4cm of gittree] (gitrepo) {
  \textbf{GitRepo}
  \nodepart{two}
  project\_name: str\\
  project\_owner\_name: str\\
  gitprovider: GitProvider\\
  access\_protocol: AccessProtocol\\
  commit\_sha: str | None
  \nodepart{three}
  resolved\_group\_name: str
};

% --- GitRunner ---
\node[classbox, right=2.4cm of client] (runner) {
  \textbf{GitRunner}
  \nodepart{two}
  executable: str
  \nodepart{three}
  clone(url, dest, branch)\\
  rev\_parse\_head(path)\\
  current\_branch(path)\\
  local\_branch\_exists(path, b)\\
  create\_branch(path, b)\\
  checkout(path, b)\\
  has\_unresolved\_merge(path)\\
  branch\_tracking\_state(path)\\
  has\_staged\_changes(path)\\
  stage\_all(path)\\
  commit(path, message)\\
  push(path, remote, branch)\\
  add\_submodule(path, url, rel, branch)
};

% --- RepoAddress ---
\node[classbox, below=1.8cm of gitrepo] (repoaddr) {
  \textbf{RepoAddress}
  \nodepart{two}
  gitprovider: GitProvider\\
  project\_name: str\\
  project\_owner\_name: str\\
  group\_name: str | None\\
  gitprovider\_url: str | None
  \nodepart{three}
  to\_ssh()\\
  to\_https()\\
  to\_url(protocol)\\
  from\_repo(repo)
};

% --- Arrows ---
% Client -> Orchestre (composition)
\draw[comporr] (client.south) -- node[right, font=\scriptsize]{1} (orchestre.north);

% Client -> GitRunner (has)
\draw[hasarrow] (client.east) -- (runner.west);

% Orchestre -> GitTree (composition)
\draw[comporr] (orchestre.south) -- node[right, font=\scriptsize]{1} (gittree.north);

% GitTree -> GitRepo (composition / dict)
\draw[comporr] (gittree.east) -- node[above, font=\scriptsize]{0..*} (gitrepo.west);

% RepoAddress -> GitRepo (uses)
\draw[hasarrow, dashed] (repoaddr.north) -- node[right, font=\scriptsize]{uses} (gitrepo.south);

\end{tikzpicture}
\end{figure}

% ============================================================
%  Figure 2 — GitTree lifecycle states
% ============================================================

\begin{figure}[ht]
\centering
\caption{GitTree lifecycle state machine.  The simplified user-facing lifecycle
  is \texttt{initialise(.cgs)} $\rightarrow$ clone $\rightarrow$ \texttt{READY};
  \texttt{initialise(.gts)} $\rightarrow$ restore $\rightarrow$ \texttt{READY};
  \texttt{pull(.cgs|.gts)} $\rightarrow$ resync $\rightarrow$ \texttt{READY}.
  Synchronized Git operations (\texttt{checkout},
  \texttt{add}, \texttt{git("commit")}, \texttt{git("push")},
  \texttt{git("tag")},
  \texttt{freeze}) are gated on \texttt{READY}.}
\label{fig:gittree-lifecycle}

\begin{tikzpicture}[
  node distance=1.4cm and 2.0cm,
  state/.style={rectangle, rounded corners=4pt, draw=blue!50, fill=blue!8,
                font=\ttfamily\small, minimum width=2.6cm, minimum height=0.8cm,
                inner sep=4pt},
  sidestate/.style={rectangle, rounded corners=4pt, draw=orange!60, fill=orange!10,
                    font=\ttfamily\small, minimum width=2.6cm, minimum height=0.8cm,
                    inner sep=4pt},
  trans/.style={-{Stealth[length=6pt]}, thick},
]

\node[state] (unloaded) {UNLOADED};
\node[state, right=of unloaded] (declared)  {DECLARED};
\node[state, right=of declared] (pending)   {PENDING};
\node[state, right=of pending]  (ready)     {READY};

\node[sidestate, below=1.2cm of declared] (partial) {PARTIAL};
\node[sidestate, below=1.2cm of pending]  (error)   {ERROR};

\draw[trans] (unloaded) -- node[above, font=\scriptsize]{load(.cgs)} (declared);
\draw[trans] (declared) -- node[above, font=\scriptsize]{discover nested} (pending);
\draw[trans] (pending)  -- node[above, font=\scriptsize]{clone/validate} (ready);
\draw[trans, dashed] (declared.south) -- (partial.north);
\draw[trans, dashed] (pending.south)  -- (error.north);
\draw[trans, bend right=20] (ready.north) to node[above, font=\scriptsize]{freeze / next .gts id + .lgr update} (declared.north);

\end{tikzpicture}
\end{figure}

% ============================================================
%  Figure 3 — Document class hierarchy
% ============================================================

\begin{figure}[ht]
\centering
\caption{Document class hierarchy.  \texttt{ConfigDocument} is the abstract
  base; each concrete subclass parses and validates one file format.}
\label{fig:document-hierarchy}

\begin{tikzpicture}[node distance=1.2cm and 1.8cm]

\node[classbox] (base) {
  \textbf{ConfigDocument}
  \nodepart{two}
  config_document.py\\
  \_data: dict
  \nodepart{three}
  read(key) / get(key)\\
  validate()\\
  from\_toml / from\_json / from\_yaml\\
  to\_toml / to\_json / to\_yaml
};

\node[classbox, below left=1.6cm and 0.6cm of base] (cgs) {
  \textbf{CgsDocument}
  \nodepart{two}
  cgs.py\\
  DOCUMENT\_KIND = ``cgs''
  \nodepart{three}
  project\_name\\
  default\_branch\\
  repos\\
  normalize authoring syntax\\
  runtime\_setting(key)
};

\node[classbox, below=1.6cm of base] (gts) {
  \textbf{GtsDocument}
  \nodepart{two}
  DOCUMENT\_KIND = ``gts''
  \nodepart{three}
  schema\_version\\
  snapshot\_hash\\
  compute\_snapshot\_hash()\\
  ensure\_snapshot\_hash()\\
  lifecycle\_state\\
  is\_ready\\
  repo\_states
};

\node[classbox, below right=1.6cm and 0.6cm of base] (goc) {
  \textbf{GocDocument}
  \nodepart{two}
  DOCUMENT\_KIND = ``goc''
  \nodepart{three}
  project\_source\\
  interaction / profile / transport\\
  actions\\
  session\_setting(key)
};

\draw[inharrow] (cgs.north) -- (base.south west);
\draw[inharrow] (gts.north) -- (base.south);
\draw[inharrow] (goc.north) -- (base.south east);

\end{tikzpicture}
\end{figure}

% ============================================================
%  Figure 4 — Registry model
% ============================================================

\begin{figure}[ht]
\centering
\caption{Registry model.  \texttt{WorkingGitTree} is the authoritative
  in-memory graph; \texttt{WorkingRepo} records both static identity
  (from \texttt{.cgs}) and dynamic state (from synchronisation operations).}
\label{fig:registry-model}

\begin{tikzpicture}[node distance=1.2cm and 2.0cm]

% WorkingGitTree
\node[classbox] (registry) {
  \textbf{WorkingGitTree}
  \nodepart{two}
  repos: dict[str, WorkingRepo]\\
  lifecycle\_state: TreeLifecycleState
  \nodepart{three}
  add(entry) / get(id) / values()\\
  children\_of(parent\_id)\\
  is\_complete() / is\_ready()\\
  recompute\_tree\_state()
};

% WorkingRepo
\node[classbox, right=2.2cm of registry] (entry) {
  \textbf{WorkingRepo}
  \nodepart{two}
  repo\_id / name / node\_type\\
  parent\_id / absolute\_path\\
  current\_ref\_kind / \_name\\
  target\_ref\_kind / \_name\\
  resolved\_ref\_kind / \_name\\
  repo\_lifecycle\_state\\
  sync\_state / discovery\_state\\
  commit\_sha / fallback\_*
  \nodepart{three}
  (mutable dataclass)
};

% ProjectTreeState
\node[classbox, below=1.6cm of registry] (treestate) {
  \textbf{ProjectTreeState}
  \nodepart{two}
  lifecycle\_state: TreeLifecycleState\\
  is\_ready: bool\\
  registry\_complete: bool
  \nodepart{three}
  (frozen dataclass)
};

% Enum: NodeType
\node[enumbox, below left=1.6cm and 0.2cm of entry] (nodetype) {
  \textbf{NodeType} \textit{(enum)}
  \nodepart{two}
  ROOT / PARENT / LEAF
};

% Enum: TreeLifecycleState
\node[enumbox, right=1.8cm of treestate] (tlc) {
  \textbf{TreeLifecycleState}
  \nodepart{two}
  UNLOADED / DECLARED\\
  PENDING / READY\\
  PARTIAL / ERROR
};

% Enum: RepoLifecycleState
\node[enumbox, below=1.2cm of nodetype] (rlc) {
  \textbf{RepoLifecycleState}
  \nodepart{two}
  DECLARED / PENDING\\
  READY / FALLBACK\_READY\\
  MISSING / ERROR
};

% Arrows
\draw[comporr] (registry.east) -- node[above, font=\scriptsize]{0..*} (entry.west);
\draw[hasarrow] (registry.south) -- (treestate.north);
\draw[hasarrow] (entry.south west) -- (nodetype.north east);
\draw[hasarrow] (entry.south) -- (rlc.north);
\draw[hasarrow] (registry.south east) -- (tlc.north west);

\end{tikzpicture}
\end{figure}


% -----------------------------------------------------------------------
\section{GitTree Lifecycle}

Figure~\ref{fig:gittree-lifecycle} shows the valid lifecycle transitions for
the \texttt{GitTree} (tracked via \texttt{DependencyTreeRegistry.lifecycle\_state}).

Valid states and their transitions:

\begin{itemize}
  \item \textbf{UNLOADED} — initial state; no data loaded.
  \item \textbf{DECLARED} — a \texttt{.cgs} file has been parsed; all entries
    exist in the registry with \texttt{DECLARED} repo lifecycle state.
  \item \textbf{PENDING} — at least one entry is awaiting synchronisation.
  \item \textbf{READY} — every entry is \texttt{READY} or
    \texttt{FALLBACK\_READY}, all commits are resolved, and the registry is
    complete.
  \item \textbf{PARTIAL} — at least one entry has not advanced past
    \texttt{DECLARED} and others have.
  \item \textbf{ERROR} — at least one entry is in the \texttt{ERROR} state.
\end{itemize}

\textbf{Hard invariants:}
\begin{itemize}
  \item \texttt{READY} means the registry is complete and synchronised — not
    that worktrees are clean.
  \item \texttt{commit}, \texttt{push}, \texttt{tag}, and
    \texttt{freeze\_release} refuse to run on a non-\texttt{READY} tree.
  \item \texttt{clone}, \texttt{restart}, \texttt{checkout}, and
    \texttt{launch\_release} must produce a \texttt{READY} tree or fail
    explicitly.
\end{itemize}

% -----------------------------------------------------------------------
\section{Document Hierarchy}

Figure~\ref{fig:document-hierarchy} shows the document class hierarchy.

All persistent documents share the \texttt{ConfigDocument} base class, which
provides format-agnostic serialisation (\texttt{to\_toml}, \texttt{to\_json},
\texttt{to\_yaml}) and factory methods (\texttt{from\_toml}, \texttt{from\_json},
\texttt{from\_yaml}).  Subclasses add document-specific validation and
convenience properties.

\begin{itemize}
  \item \textbf{CgsDocument} (\texttt{.cgs}) — the authoring spec; describes
    the static project topology.  \emph{Never} a runtime snapshot.
  \item \textbf{GtsDocument} (\texttt{.gts}) — a generated snapshot capturing
    the exact state of the tree including commit SHAs and absolute paths.
    \emph{Never} hand-edited.
  \item \textbf{GocDocument} (\texttt{.goc}) — a command script encoding a
    sequence of \texttt{cgitsync} commands with session defaults.
\end{itemize}

% -----------------------------------------------------------------------
\section{Registry Model}

Figure~\ref{fig:registry-model} shows the registry model.

\texttt{DependencyTreeRegistry} is the authoritative in-memory graph.  It maps
\emph{repo IDs} (colon-separated path strings such as
\texttt{root:deps/child-repo}) to \texttt{RepoRegistryEntry} instances.
\texttt{RepoRegistryEntry} holds both the static identity declared in
\texttt{.cgs} and the dynamic state updated during operations.

Builder functions in \texttt{registry.py} translate document objects into a
live registry:

\begin{lstlisting}[language=Python]
# Load from authoring spec
registry = build_registry_from_cgs_document(document, config_path)

# Load from snapshot
registry = build_registry_from_gts_document(snapshot_document)

# Serialise to snapshot
gts_doc = build_gts_document_from_registry(
    registry, command_origin="clone", source_cgs_path=path
)
\end{lstlisting}
